\begin{table}
\centering
\begin{tabular}{llll|lll|lll}
\toprule
     & \multicolumn{3}{c}{Try Again ($n$ = 10)} & \multicolumn{3}{c}{ReAct ($n$ = 10)} & \multicolumn{3}{c}{Plan \& Solve} \\
     & SR & Turns & Error \%
     & SR & Turns & Error \%
     & SR & Turns & Error \% \\
\midrule
SQL & 47.3 & 7.25 & 46.4
    & \textbf{58.7} & 5.30 & \textbf{6.94}
    & 49.1 & \textbf{4.29} & 16.2 \\
Bash & \textbf{46.5} & 6.15 & 24.9
    & 20.5 & \textbf{4.40} & \textbf{20.4}
    & 28.0 & 6.65 & 53.3 \\
\bottomrule
\end{tabular}
\vspace{5pt}
\caption{Comparison of different prompting strategies across the entire \benchmark{}-SQL and \benchmark{}-Bash datasets using \texttt{gpt-3.5-turbo} as the base model. \textit{Turns} refers to the average number of turns taken for a single task episode. For Try Again and ReAct, the max number of turns $n = 10$. The highest Success Rate, fewest Turns, and lowest Error \% are highlighted per dataset since they reflect more accuracy and efficient task solving. Best metrics are in \textbf{bold}.}
\label{table:strategy-comparison}
\vspace{-1.5em}
\end{table}